\chapter{Fase 5: Conectar el teléfono real}

\section{Objetivo}

Reemplazar el softphone por un teléfono analógico DTMF físico, conectado a la centralita a través del adaptador Grandstream HT801. Al terminar esta fase, descolgar el teléfono debe llevar directamente al juego, sin marcar ningún número.

\section{Conceptos previos}

\subsection{Qué es un puerto FXS}

Un puerto \emph{FXS (Foreign eXchange Subscriber)} es el tipo de puerto que provee la compañía telefónica al usuario doméstico: entrega voltaje continuo de aproximadamente 48\,V para alimentar el teléfono, voltaje alterno de unos 90\,V para hacerlo sonar, y mediante el \emph{loop current} detecta cuándo se descuelga o cuelga el aparato.

Un puerto \emph{FXO (Foreign eXchange Office)} es el opuesto: se conecta hacia una central telefónica y se comporta como un teléfono.

El HT801 expone un puerto FXS por su conector RJ11, lo cual permite enchufar cualquier teléfono analógico como si la central fuera la compañía telefónica.

\subsection{Auto-dial al descolgar}

La función \emph{Off-hook Auto-Dial} del HT801 permite configurar un número que se marcará automáticamente al descolgar el teléfono. Esta es la clave para replicar la experiencia Hugo: el usuario descuelga e inmediatamente está jugando, sin necesidad de marcar nada.

\section{Conexiones físicas}

\begin{center}
\begin{tikzpicture}[
  node distance=2cm,
  every node/.style={font=\small},
  box/.style={draw=hugoblue, thick, rounded corners=4pt, fill=hugoblue!10, align=center}
]
  \node[box, minimum width=2.8cm, minimum height=1.2cm] (phone) at (0,0) {Teléfono\\analógico};
  \node[box, minimum width=2.8cm, minimum height=1.2cm] (ata) at (5,0) {HT801};
  \node[box, minimum width=2.8cm, minimum height=1.2cm] (router) at (5,-2.5) {Router LAN};
  \node[box, minimum width=2.8cm, minimum height=1.2cm] (pi) at (10,-2.5) {Raspberry Pi};

  \draw[->, thick, hugoaccent] (phone) -- node[above,font=\tiny] {RJ11 (PHONE)} (ata);
  \draw[->, thick, hugoaccent] (ata) -- node[right,font=\tiny] {RJ45 (LAN)} (router);
  \draw[<-, thick, hugoaccent] (router) -- (pi);

  \node[below=0.3cm of ata, font=\tiny, color=gray] {alim. 5V};
  \node[below=0.3cm of phone, font=\tiny, color=gray] {sin alim. externa};
\end{tikzpicture}
\end{center}

\begin{itemize}
  \item El cable RJ11 del teléfono va al puerto rotulado \textbf{PHONE} (o \textbf{FXS}) del HT801. \emph{No} al puerto LAN.
  \item El HT801 va por Ethernet al mismo switch/router donde está la Pi.
  \item La fuente del HT801 (5\,V) se conecta al puerto correspondiente.
  \item El teléfono analógico no necesita alimentación externa: la línea le provee voltaje.
\end{itemize}

\section{Paso 1: Descubrir la IP del HT801}

El HT801 obtiene IP por DHCP automáticamente. Para conocerla hay dos métodos:

\subsection{Método A: IVR del propio HT801}

\begin{enumerate}
  \item Descolgar el teléfono conectado al HT801.
  \item Marcar \texttt{***} (tres asteriscos).
  \item Una voz en inglés ofrece un menú de opciones de diagnóstico.
  \item Marcar \texttt{02} para escuchar la IP del dispositivo.
\end{enumerate}

\subsection{Método B: panel del router}

Entrar al panel de administración del router y buscar en la lista de clientes DHCP un dispositivo nuevo con nombre similar a \texttt{GS\_HT801} o \texttt{HT801-xxxx}.

Una vez identificada la IP (supongamos \texttt{192.168.1.11}), \textbf{reservar} esa IP para la MAC del HT801 en el router, igual que se hizo con la Pi.

\section{Paso 2: Acceder al panel web del HT801}

Desde la Mac, abrir en el navegador:

\begin{plaincode}
http://192.168.1.11
\end{plaincode}

Credenciales:

\begin{itemize}
  \item \textbf{Usuario:} \texttt{admin}
  \item \textbf{Contraseña:} \texttt{admin} (modelos antiguos) o la contraseña impresa en la etiqueta inferior del dispositivo (modelos recientes con contraseña única por unidad).
\end{itemize}

\begin{advertencia}
Cambiar la contraseña por defecto antes de hacer cualquier otra cosa: \textsf{Advanced Settings} $\to$ \textsf{Admin Password}.
\end{advertencia}

\section{Paso 3: Configurar la cuenta SIP}

Navegar a la pestaña \textsf{FXS Port} (o \textsf{Profile 1} en algunos firmwares).

\begin{table}[h]
\centering
\begin{tabular}{ll}
\toprule
\textbf{Campo} & \textbf{Valor} \\
\midrule
Account Active & Yes \\
Primary SIP Server & 192.168.1.10 (IP de la Pi) \\
SIP User ID & 101 \\
Authenticate ID & 101 \\
Authenticate Password & el secreto de 101 en \texttt{sip.conf} \\
Name & Telefono Hugo 1 \\
DTMF Payload Type & 101 \\
Preferred DTMF method & RFC2833 (priority 1) \\
SIP Registration & Yes \\
Outgoing call without registration & No \\
NAT Traversal & No \\
\bottomrule
\end{tabular}
\caption{Configuración SIP del HT801.}
\end{table}

En la sección de codecs:

\begin{itemize}
  \item Choice 1: \texttt{PCMU} (G.711 $\mu$-law)
  \item Choice 2: \texttt{PCMA} (G.711 A-law)
  \item Resto: \texttt{None}
\end{itemize}

\section{Paso 4: Configurar Auto-Dial}

En la misma pestaña, buscar el campo \texttt{Offhook Auto-Dial} y configurarlo en:

\begin{plaincode}
9000
\end{plaincode}

Esto provoca que al descolgar el teléfono, el HT801 marque automáticamente al \texttt{9000}, que es la extensión que el dialplan asocia con el juego.

Hacer clic en \textsf{Update} y luego en \textsf{Apply}. El HT801 reinicia su módulo SIP y queda registrado al cabo de unos 10 segundos.

\section{Paso 5: Verificar registro}

En la Pi:

\begin{shellcode}
sudo asterisk -rvvv
\end{shellcode}

\begin{plaincode}
hugo-pbx*CLI> sip show peers
Name/username             Host                Status
101/101                   192.168.1.11        OK (5 ms)
102/102                   (Unspecified)       UNKNOWN
\end{plaincode}

La extensión 101 debe aparecer con la IP del HT801 y estado \texttt{OK}. Si dice \texttt{UNREACHABLE} o el campo Host aparece vacío, verificar:

\begin{itemize}
  \item Coincidencia de contraseñas entre HT801 y \texttt{sip.conf}.
  \item Que la IP de \emph{Primary SIP Server} apunte realmente a la Pi.
  \item Los logs detallados:
  \begin{shellcode}
hugo-pbx*CLI> sip set debug on
  \end{shellcode}
\end{itemize}

\section{Paso 6: La prueba final}

Todas las piezas están conectadas. La secuencia:

\begin{enumerate}
  \item Confirmar que DOSBox-Staging muestra Hugo en la TV.
  \item Confirmar que el servicio bridge está activo: \texttt{systemctl status hugo-bridge}.
  \item \textbf{Descolgar el teléfono.}
  \item En el auricular debe escucharse el beep del dialplan (después de un breve momento mientras el HT801 marca 9000 internamente).
  \item Marcar \textbf{2}. Hugo salta.
  \item Marcar \textbf{4}, \textbf{6}, \textbf{8}, \textbf{5} y jugar normalmente.
  \item Al colgar el teléfono, la llamada finaliza limpiamente.
\end{enumerate}

Si todo funciona, el sistema está completo en su configuración mínima. El proyecto Hugo Phone está vivo.

\section{Resolución de problemas}

\subsection{El teléfono no da tono al descolgar}

El HT801 no está dando línea. Causas posibles:

\begin{itemize}
  \item LED del puerto FXS en rojo: error de hardware o configuración. Hacer un \emph{factory reset} (botón Reset físico presionado 7 segundos).
  \item Cable RJ11 en mal estado o invertido (probar otro).
  \item Teléfono dañado (probar con otro teléfono).
\end{itemize}

\subsection{Hay tono pero al descolgar no marca el 9000}

\emph{Offhook Auto-Dial} no está bien configurado. Verificar que sea exactamente \texttt{9000} sin espacios, sin asteriscos, sin numerales.

\subsection{Conecta pero no se escucha el beep}

Problema de codecs:

\begin{itemize}
  \item Asegurar que el HT801 prioriza \texttt{PCMU} y \texttt{PCMA}.
  \item En el HT801, sección \emph{Audio Settings}, subir las ganancias \emph{Tx} y \emph{Rx} si el volumen es muy bajo.
\end{itemize}

\subsection{Los DTMF no producen acción en Hugo}

\begin{itemize}
  \item Verificar el método DTMF en el HT801: debe ser \texttt{RFC2833}.
  \item En \texttt{sudo asterisk -rvvv} debe aparecer el \texttt{UserEvent: HugoKey} al pulsar dígitos.
  \item Algunos HT801 antiguos fallan con RFC 2833 y requieren cambiar a \texttt{SIP INFO}. Probar ese método si el anterior no funciona.
\end{itemize}

\subsection{La llamada se cae sola después de algunos minutos}

El dialplan tiene un timeout de 3600 segundos (1 hora) en \texttt{Read()}. Si la llamada cae antes, probablemente un \emph{session timer} está activo. Verificar en \texttt{sip.conf} la línea \texttt{session-timers=refuse}.
